% ========================================================
% BÀI TẬP TOÁN 5 - KẾT NỐI TRI THỨC VỚI CUỘC SỐNG
% Bài 73: Ôn tập toán chuyển động đều
% Định dạng: BTTeX Part III - Trắc nghiệm điền số ngắn
% ========================================================

% Câu 1
\begin{ex}
  Một ô tô đi quãng đường $150\text{ km}$ hết $3$ giờ. Vận tốc của ô tô là bao nhiêu km/h?
  \shortans{50}
  \loigiai{
    $150 : 3 = 50\text{ km/h}$.
  }
\end{ex}

% Câu 2
\begin{ex}
  Một người đi xe máy với vận tốc $36\text{ km/h}$ trong $2$ giờ $30$ phút ($2{,}5$ giờ). Quãng đường đi được là bao nhiêu ki-lô-mét?
  \shortans{90}
  \loigiai{
    $36 \times 2{,}5 = 90\text{ km}$.
  }
\end{ex}

% Câu 3
\begin{ex}
  Một người đi bộ với vận tốc $4\text{ km/h}$ trên quãng đường dài $14\text{ km}$. Thời gian đi đổi ra phút là bao nhiêu phút?
  \shortans{210}
  \loigiai{
    $14 : 4 = 3{,}5\text{ giờ} = 210$ phút.
  }
\end{ex}

% Câu 4
\begin{ex}
  Hai xe ô tô cùng xuất phát từ hai điểm cách nhau $180\text{ km}$ đi ngược chiều. Vận tốc xe thứ nhất là $55\text{ km/h}$, xe thứ hai là $35\text{ km/h}$. Sau mấy giờ hai xe gặp nhau?
  \shortans{2}
  \loigiai{
    Tổng vận tốc: $55 + 35 = 90\text{ km/h}$. Thời gian gặp: $180 : 90 = 2\text{ giờ}$.
  }
\end{ex}

% Câu 5
\begin{ex}
  Ở bài trên, điểm gặp nhau cách điểm xuất phát của xe thứ nhất bao nhiêu ki-lô-mét?
  \shortans{110}
  \loigiai{
    $55 \times 2 = 110\text{ km}$.
  }
\end{ex}

% Câu 6
\begin{ex}
  Một người đi xe đạp từ $A$ lúc $7\text{ giờ}$ với vận tốc $12\text{ km/h}$. Đến $9\text{ giờ}$, một xe máy cũng từ $A$ đuổi theo với vận tốc $36\text{ km/h}$. Sau mấy giờ xe máy đuổi kịp xe đạp?
  \shortans{1}
  \loigiai{
    Quãng đường xe đạp đi trước: $12 \times 2 = 24\text{ km}$. Hiệu vận tốc: $36 - 12 = 24\text{ km/h}$. Thời gian: $24 : 24 = 1\text{ giờ}$.
  }
\end{ex}

% Câu 7
\begin{ex}
  Ở bài trên, xe máy đuổi kịp xe đạp lúc mấy giờ?
  \shortans{10}
  \loigiai{
    $9 + 1 = 10\text{ giờ}$.
  }
\end{ex}

% Câu 8
\begin{ex}
  Một ca nô có vận tốc thực $22\text{ km/h}$, vận tốc dòng nước là $3\text{ km/h}$. Vận tốc ca nô khi xuôi dòng là bao nhiêu km/h?
  \shortans{25}
  \loigiai{
    $22 + 3 = 25\text{ km/h}$.
  }
\end{ex}

% Câu 9
\begin{ex}
  Ở bài trên, vận tốc ca nô khi ngược dòng là bao nhiêu km/h?
  \shortans{19}
  \loigiai{
    $22 - 3 = 19\text{ km/h}$.
  }
\end{ex}

% Câu 10
\begin{ex}
  Một ca nô xuôi dòng khúc sông dài $50\text{ km}$ với vận tốc $25\text{ km/h}$. Thời gian ca nô đi hết bao nhiêu giờ?
  \shortans{2}
  \loigiai{
    $50 : 25 = 2\text{ giờ}$.
  }
\end{ex}

% Câu 11
\begin{ex}
  Một đoàn tàu dài $150\text{ m}$ chạy qua một cột điện hết $10$ giây. Vận tốc của đoàn tàu là bao nhiêu m/s?
  \shortans{15}
  \loigiai{
    $150 : 10 = 15\text{ m/s}$.
  }
\end{ex}

% Câu 12
\begin{ex}
  Đổi vận tốc $15\text{ m/s}$ ra km/h được bao nhiêu km/h?
  \shortans{54}
  \loigiai{
    $15 \times 3{,}6 = 54\text{ km/h}$.
  }
\end{ex}

% Câu 13
\begin{ex}
  Hai thị xã $A$ và $B$ cách nhau $100\text{ km}$. Một người đi xe máy từ $A$ đến $B$ với vận tốc $40\text{ km/h}$. Thời gian đi đổi ra phút là bao nhiêu phút?
  \shortans{150}
  \loigiai{
    $100 : 40 = 2{,}5\text{ giờ} = 150$ phút.
  }
\end{ex}

% Câu 14
\begin{ex}
  Một ô tô đi từ $A$ lúc $6\text{ giờ } 15\text{ phút}$ và đến $B$ lúc $8\text{ giờ } 45\text{ phút}$, dọc đường nghỉ $30$ phút. Thời gian ô tô thực chạy là bao nhiêu giờ?
  \shortans{2}
  \loigiai{
    Tổng thời gian: $2{,}5$ giờ. Thực chạy: $2{,}5 - 0{,}5 = 2$ giờ.
  }
\end{ex}

% Câu 15
\begin{ex}
  Nếu quãng đường ở câu trên dài $100\text{ km}$, vận tốc ô tô là bao nhiêu km/h?
  \shortans{50}
  \loigiai{
    $100 : 2 = 50\text{ km/h}$.
  }
\end{ex}

% Câu 16
\begin{ex}
  Một máy bay bay với vận tốc $750\text{ km/h}$ trong $1$ giờ $48$ phút ($1{,}8$ giờ). Quãng đường bay được là bao nhiêu ki-lô-mét?
  \shortans{1350}
  \loigiai{
    $750 \times 1{,}8 = 1\,350\text{ km}$.
  }
\end{ex}

% Câu 17
\begin{ex}
  Một người đi xe đạp với vận tốc $12\text{ km/h}$. Để đi quãng đường $18\text{ km}$ người đó cần bao nhiêu phút?
  \shortans{90}
  \loigiai{
    $18 : 12 = 1{,}5\text{ giờ} = 90$ phút.
  }
\end{ex}

% Câu 18
\begin{ex}
  Một xe buýt đi quãng đường $30\text{ km}$ hết $45$ phút ($0{,}75$ giờ). Vận tốc của xe buýt là bao nhiêu km/h?
  \shortans{40}
  \loigiai{
    $30 : 0{,}75 = 40\text{ km/h}$.
  }
\end{ex}

% Câu 19
\begin{ex}
  Hai xe đạp cùng xuất phát một lúc đi cùng chiều: xe thứ nhất đi $14\text{ km/h}$, xe thứ hai đi $10\text{ km/h}$. Sau $3$ giờ khoảng cách giữa hai xe là bao nhiêu ki-lô-mét?
  \shortans{12}
  \loigiai{
    $(14 - 10) \times 3 = 12\text{ km}$.
  }
\end{ex}

% Câu 20
\begin{ex}
  Một vận động viên bơi $50\text{ m}$ hết $25$ giây. Vận tốc bơi của vận động viên là bao nhiêu m/s?
  \shortans{2}
  \loigiai{
    $50 : 25 = 2\text{ m/s}$.
  }
\end{ex}

% Câu 21
\begin{ex}
  Một ô tô chạy từ $A$ đến $B$ dài $140\text{ km}$ với vận tốc $56\text{ km/h}$. Thời gian ô tô chạy là bao nhiêu phút?
  \shortans{150}
  \loigiai{
    $140 : 56 = 2{,}5\text{ giờ} = 150$ phút.
  }
\end{ex}

% Câu 22
\begin{ex}
  Một người bơi xuôi dòng được $800\text{ m}$ trong $10$ phút. Vận tốc bơi xuôi dòng là bao nhiêu m/phút?
  \shortans{80}
  \loigiai{
    $800 : 10 = 80\text{ m/phút}$.
  }
\end{ex}

% Câu 23
\begin{ex}
  Quãng đường từ nhà An đến trường dài $1\,500\text{ m}$. An đi bộ với vận tốc $75\text{ m/phút}$. An đi đến trường hết bao nhiêu phút?
  \shortans{20}
  \loigiai{
    $1\,500 : 75 = 20$ phút.
  }
\end{ex}

% Câu 24
\begin{ex}
  Một người đi xe máy từ nhà về quê hết $1\text{ giờ } 45\text{ phút}$ với vận tốc $40\text{ km/h}$. Quãng đường dài bao nhiêu ki-lô-mét?
  \shortans{70}
  \loigiai{
    $40 \times 1{,}75 = 70\text{ km}$.
  }
\end{ex}

% Câu 25
\begin{ex}
  Hai người đi bộ ngược chiều nhau trên một con đường dài $9\text{ km}$. Người thứ nhất đi $5\text{ km/h}$, người thứ hai đi $4\text{ km/h}$. Sau mấy giờ hai người gặp nhau?
  \shortans{1}
  \loigiai{
    $9 : (5 + 4) = 1\text{ giờ}$.
  }
\end{ex}
